% Circuito de medición para el Experimento 2: Tren de pulsos con DSO y Multímetro True RMS
\begin{circuitikz}[scale=0.9, transform shape]
    % Generador de funciones
    \draw (0,0.5) node[ground] {}
        to[sV, l={$G_1$ (Pulsos)}] (0,2.5)
        to[short, -*] (2.2,2.5);
    \draw (0,0.5) to[short, -*] (2.2,0.5);

    % Derivación superior hacia el Osciloscopio DSO
    \draw (2.2,2.5) to[short] (2.2,2.2)
        to[short, -o] (4.8,2.2) node[above, font=\scriptsize\bfseries] {CH1};
    \draw (2.2,0.5) to[short] (2.2,1.0)
        to[short, -o] (4.8,1.0) node[above, font=\scriptsize\bfseries] {Masa};

    % Conexión al bloque DSO
    \draw (4.8,2.2) to[short] (5.4,2.2);
    \draw (4.8,1.0) to[short] (5.4,1.0);
    \draw (4.8,1.0) node[ground] {};

    % Bloque osciloscopio
    \draw[thick, rounded corners=3pt, fill=blue!3] (5.4,0.4) rectangle (10.0,2.9);
    \node at (7.7,2.5) [font=\bfseries\small] {Osciloscopio DSO};
    \node at (7.7,1.85) [font=\scriptsize, align=center] {CH1: FFT / Temporal\\Promedio 64 cuentas};
    \node at (7.7,1.1) [font=\scriptsize, align=center] {Ventanas: Flattop / Rect.\\Cursores de frecuencia y amplitud};

    % Derivación inferior hacia el Multímetro True RMS
    \draw (2.2,2.5) to[short] (2.2,-1.1)
        to[short, -o] (4.8,-1.1) node[above, font=\scriptsize\bfseries] {$V/\Omega$};
    \draw (2.2,0.5) to[short] (2.2,-1.9)
        to[short, -o] (4.8,-1.9) node[above, font=\scriptsize\bfseries] {COM};

    % Conexión al bloque Multímetro
    \draw (4.8,-1.1) to[short] (5.4,-1.1);
    \draw (4.8,-1.9) to[short] (5.4,-1.9);
    \draw (4.8,-1.9) node[ground] {};

    % Bloque Multímetro True RMS
    \draw[thick, rounded corners=3pt, fill=purple!5] (5.4,-2.5) rectangle (10.0,-0.6);
    \node at (7.7,-1.1) [font=\bfseries\small] {Multímetro Digital};
    \node at (7.7,-1.8) [font=\scriptsize, align=center] {UNI-T UT89X (Escala $V_{\text{DC}}$)\\Medición directa del valor medio};
\end{circuitikz}
